% =====================================================================
\section{Especificación de requisitos final}
\label{sec:ers}

\subsection{Catálogo final de requisitos}

\begin{table}[H]
\centering\footnotesize
\caption{Composición del catálogo de requisitos del ERS de SGCV-IA (v3.0)}
\label{tab:catalogo}
\begin{tabular}{|L{5.6cm}|C{2.0cm}|C{2.4cm}|L{4.4cm}|}
\hline
\rowcolor{grisclaro}\textbf{Tipo} & \textbf{Cantidad} & \textbf{Rango} & \textbf{Observación}\\ \hline
Requisitos funcionales & 25 & \id{RF-01}--\id{RF-24}, \id{RF-26} & 19 \emph{Must}, 6 \emph{Should}; \id{RF-25} reservado y reasignado por el CCB\\ \hline
Requisitos no funcionales & 18 & \id{RNF-01}--\id{RNF-18} & Mapeados a características de ISO/IEC 25010; supera el mínimo de 15 exigido por la rúbrica\\ \hline
Requisitos legales & 14 & \id{RL-01}--\id{RL-14} & Derivados de la LOPDP, trazados hacia RF/RNF\\ \hline
Restricciones de diseño & 9 & \id{RST-01}--\id{RST-09} & Presupuestaria, tecnológica, técnica/hardware, arquitectónica, ética/regulatoria y de seguridad\\ \hline
\textbf{Total formalizado} & \textbf{66} & --- & A partir de 24 requisitos crudos (\id{RC-01}--\id{RC-24}) elicitados en la Unidad II\\ \hline
\end{tabular}
\end{table}

\noindent
El catálogo de 25 RF se formalizó a partir de veinticuatro requisitos crudos elicitados
(\id{RC-01} a \id{RC-24}), más dos requisitos legales derivados directamente de la LOPDP
(\id{RL-01} y el bloque \id{RL-04}/\id{RL-05}/\id{RL-06}/\id{RL-07}, Criterio C4). Uno de los
crudos, \id{RC-22} (``Bajo costo de implementación y operación''), se reclasificó como
restricción de diseño (\id{RST-01}) por tratarse de un criterio de negocio y no de una función
observable del sistema, lo que explica por qué el rango funcional salta de \id{RF-24} a
\id{RF-26}: \id{RF-25} fue aprobado por el Change Control Board en una práctica anterior para un
requisito de contenido distinto (asistencia a citas) y el identificador original quedó reservado
para evitar colisión.

\subsection{Requisitos clave y por qué lo son}

De los veinticinco requisitos funcionales, cinco concentran la propuesta de valor del sistema y
son los que conviene poder defender uno por uno.

\begin{table}[H]
\centering\footnotesize
\caption{Requisitos clave de SGCV-IA y su justificación en la evidencia}
\label{tab:req-clave}
\begin{tabular}{|C{1.4cm}|L{4.2cm}|L{4.4cm}|L{4.4cm}|}
\hline
\rowcolor{grisclaro}\textbf{ID} & \textbf{Requisito} & \textbf{Evidencia que lo sustenta} & \textbf{Por qué es clave}\\ \hline
\id{RF-04} & Registro detallado del historial clínico & \id{EV-01}, \id{EV-02}: entrevistas al Dr. Antonio y a Amagua & Es la base de la que dependen el diagnóstico asistido, el seguimiento nutricional y la trazabilidad de auditoría\\ \hline
\id{RF-17} & Sugerencias diagnósticas asistidas por IA & \id{EV-01}, \id{EV-02} & Es la funcionalidad distintiva del sistema frente a un gestor clínico convencional, y la que activa la mayor parte del marco normativo (\id{RL-09}, \id{RL-10}, \id{RST-05}, \id{RST-08})\\ \hline
\id{RF-21} & Operación con baja dependencia de conectividad permanente & \id{EV-01} & Sin sincronización diferida, el sistema queda inoperable en las condiciones reales de conectividad de la clínica\\ \hline
\id{RF-24} & Gestión de consentimiento informado (LOPDP) & \id{EV-04}: consentimiento informado firmado & Es la puerta de entrada legal para tratar cualquier dato del propietario o de la mascota; sin él, ningún otro RF es lícito\\ \hline
\id{RF-26} & Portal de ejercicio de derechos (acceso, rectificación, eliminación) & \id{EV-04} & Concentra la implementación de cuatro requisitos legales (\id{RL-04} a \id{RL-07}) en un único punto de contacto para el titular\\ \hline
\end{tabular}
\end{table}

\subsection{Priorización MoSCoW del catálogo}

\begin{table}[H]
\centering\footnotesize
\caption{Priorización de requisitos mediante MoSCoW}
\label{tab:moscow}
\begin{tabular}{|L{2.0cm}|L{6.0cm}|L{4.2cm}|L{2.2cm}|}
\hline
\rowcolor{grisclaro}\textbf{Prioridad} & \textbf{RF} & \textbf{RNF} & \textbf{\%}\\ \hline
Must have & \id{RF-01} a \id{RF-17}, \id{RF-24}, \id{RF-26} (19 RF) & \id{RNF-01}, \id{RNF-02}, \id{RNF-04}, \id{RNF-05}, \id{RNF-09}, \id{RNF-16} (6 RNF) & 71\,\%\\ \hline
Should have & \id{RF-18} a \id{RF-23} (6 RF) & \id{RNF-03}, \id{RNF-06}, \id{RNF-07} (3 RNF) & 25\,\%\\ \hline
Could have & Ninguno & \id{RNF-08}, \id{RNF-10} (2 RNF) & 4\,\%\\ \hline
Won't have & --- & --- & 0\,\%\\ \hline
\end{tabular}
\end{table}

\noindent
Los diecinueve \emph{Must} son las funciones sin las cuales el sistema no reemplaza el proceso
manual actual de la clínica: registro de pacientes, historial clínico, atención médica con
diagnóstico asistido por IA, control de inventario, gestión de empleados, facturación y
comunicación. Los seis \emph{Should} aportan valor significativo ---reportes comparativos,
seguimiento nutricional detallado, recordatorios automáticos, trazabilidad de auditoría--- pero el
sistema es funcional sin ellos en una primera versión. Ningún RF quedó en \emph{Could have}; a
nivel de RNF, la portabilidad extendida a más plataformas móviles (\id{RNF-08}) y la resiliencia
total ante una caída prolongada del proveedor externo de IA (\id{RNF-10}) son mejoras deseables que
no bloquean la operación básica y se abordan en una segunda iteración. No se identificaron
requisitos a excluir explícitamente del alcance funcional ya definido, porque las exclusiones de
alcance ---diagnósticos definitivos sin validación humana, prescripción automatizada, integración
con sistemas gubernamentales de salud animal, facturación electrónica certificada ante el SRI,
\emph{e-commerce}--- ya se filtraron en la fase de planificación (Entrega 1A) y nunca entraron como
candidatos a RC o RF.

\subsection{Requisitos no funcionales por característica de calidad}

Los dieciocho requisitos no funcionales se organizan según las características de calidad de
producto de ISO/IEC 25010 \cite{iso25010}. Cubren desempeño (tiempo de respuesta,
comportamiento temporal, utilización de recursos), fiabilidad (madurez, disponibilidad,
recuperabilidad), seguridad (control de acceso, trazabilidad y no repudio, cifrado de datos),
usabilidad (capacidad de aprendizaje, satisfacción del usuario), compatibilidad, mantenibilidad
(modularidad, capacidad de prueba) y dos características añadidas por el propio dominio del
proyecto: explicabilidad de la IA (\id{RNF-16}) y privacidad, esta última desdoblada en
conservación de datos (\id{RNF-17}) y respuesta a incidentes (\id{RNF-18}), ambas trazadas
directamente a artículos de la LOPDP. \id{RNF-16} es el requisito que cierra el piso de
explicabilidad exigido por el componente de IA: ninguna sugerencia diagnóstica puede presentarse
sin su referencia bibliográfica y su nivel de confianza asociado.

\subsection{Requisitos legales y su trazabilidad hacia RF/RNF}

\begin{table}[H]
\centering\footnotesize
\caption{Selección de requisitos legales derivados de la LOPDP y su implementación}
\label{tab:rl}
\begin{tabular}{|C{1.4cm}|C{1.6cm}|L{5.6cm}|L{4.2cm}|}
\hline
\rowcolor{grisclaro}\textbf{ID} & \textbf{Artículo} & \textbf{Obligación} & \textbf{Requisito que la implementa}\\ \hline
\id{RL-01} & Art. 7, 8 & Captura de consentimiento informado del propietario & \id{RF-24}\\ \hline
\id{RL-04} & Art. 12 & Informar al titular sobre las finalidades del tratamiento & \id{RF-26}\\ \hline
\id{RL-05} & Art. 13 & Derecho de acceso, atendible en 15 días & \id{RF-26}\\ \hline
\id{RL-06} & Art. 14 & Derecho de rectificación & \id{RF-26}\\ \hline
\id{RL-07} & Art. 15 & Derecho de eliminación, atendible en 15 días & \id{RF-26}\\ \hline
\id{RL-09} & Art. 20 & Ninguna sugerencia de IA se aplica sin revisión del veterinario & \id{RF-17}, \id{RF-18}, \id{RF-19}\\ \hline
\id{RL-10} & Art. 25, 26, 30, 31 & Datos de salud como categoría especial bajo secreto profesional & \id{RF-17}, \id{RNF-09}, \id{RNF-12}\\ \hline
\id{RL-13} & Art. 43, 46 & Notificación de incidentes a la Autoridad ($\leq$5 días) y al titular ($\leq$3 días) & \id{RNF-18}, \id{RF-23}\\ \hline
\end{tabular}
\end{table}

\noindent
Se listan ocho de los catorce requisitos legales, los que condicionan directamente un requisito
funcional o no funcional concreto; el catálogo completo (\id{RL-01} a \id{RL-14}) reside en la
Sección 3.4 del ERS. Los cuatro derechos del titular ---acceso, rectificación, eliminación y
portabilidad--- se concentran en un único requisito funcional (\id{RF-26}, portal de ejercicio de
derechos) en lugar de dispersarse en pantallas independientes, decisión que simplifica tanto la
implementación como la auditoría de cumplimiento.

\subsection{Restricciones de diseño}

Se identificaron nueve restricciones de diseño (\id{RST-01} a \id{RST-09}), agrupadas en cinco
tipos: presupuestaria, tecnológica, técnica/hardware, arquitectónica y ética/regulatoria-seguridad.
Ninguna proviene de una preferencia técnica del equipo: todas se derivan del alcance definido en la
Entrega 1A o de un requisito funcional o no funcional ya formalizado.

\begin{table}[H]
\centering\footnotesize
\caption{Restricciones de diseño de SGCV-IA (selección)}
\label{tab:rst}
\begin{tabular}{|C{1.4cm}|L{6.8cm}|L{2.4cm}|L{2.4cm}|}
\hline
\rowcolor{grisclaro}\textbf{ID} & \textbf{Restricción} & \textbf{Tipo / Origen} & \textbf{Relacionado}\\ \hline
\id{RST-01} & Modelo de costos accesible para clínicas de 1 a 3 veterinarios; se prioriza software de bajo costo o código abierto & Presupuestaria & \id{RC-22}\\ \hline
\id{RST-03} & Sincronización offline mediante cola local de operaciones pendientes, sin réplica completa de la base de datos en el cliente & Técnica / Hardware & \id{RF-21}\\ \hline
\id{RST-04} & El módulo de IA se implementa desacoplado del backend transaccional, de modo que su indisponibilidad no afecte al resto del sistema & Arquitectónica & \id{RF-17}, \id{RF-18}, \id{RF-19}\\ \hline
\id{RST-05} & Toda sugerencia de IA queda sujeta a revisión y aprobación explícita del veterinario antes de incorporarse al historial & Ética / Regulatoria & \id{RF-17}\\ \hline
\id{RST-08} & El sistema no emite diagnósticos definitivos ni prescripciones automatizadas; toda salida de IA es sugerencia & Regulatoria / Alcance & \id{RF-17}, \id{RF-19}\\ \hline
\end{tabular}
\end{table}

\subsection{Matriz de trazabilidad}

La matriz de trazabilidad de la \S\ref{sec:gestion} conecta evidencia de campo, requisito y caso de
uso. Las evidencias \id{EV-01} a \id{EV-04} corresponden al registro del Apéndice B: \id{EV-01} es
la entrevista al Dr. Antonio (24/05/2026); \id{EV-02}, la entrevista a Amagua (25/05/2026);
\id{EV-03}, el acta de reunión firmada; \id{EV-04}, el consentimiento informado firmado. Cuando un
requisito no corresponde a uno de los diez casos de uso detallados, se referencia el caso de uso
equivalente del diagrama general. La numeración de esta cadena de trazabilidad se mantiene
coherente con la reclasificación de \id{RC-22} señalada en la \S\ref{sec:ers}: al no ocupar fila de
RF, la tabla no presenta un salto sin explicación entre \id{RF-24} y \id{RF-26}.
